% !TeX root = ../thesis.tex
\chapter{Literature Review}
\label{ch:literature-review}

\section{Introduction}

\subsection{Scope and Structure of the Review}
\topics{whole-pipeline coverage from radiotracer to clinic; the cost lens applied
throughout; what is deliberately out of scope (e.g.\ SPECT, radiochemistry depth);
how the sections build toward the research gap.}

\subsection{Review Methodology}
\topics{databases and search engines used (PubMed, IEEE Xplore, Scopus, Google Scholar);
search terms and inclusion/exclusion criteria; citation snowballing from landmark
reviews; date ranges; how grey literature (vendor white papers, conference records) is
treated.}

\section{Historical Development of Positron Emission Tomography}

\subsection{Origins: Positron Imaging Before Tomography}
\topics{positron discovery (Anderson, 1932); first medical use of positron emitters;
early coincidence brain probes (Wrenn 1951, Brownell \& Sweet 1953); rectilinear
scanners and the pre-tomographic era.}

\subsection{The First Tomographic Scanners}
\topics{PC-I and PETT series (Ter-Pogossian, Phelps, Hoffman, 1970s); analytical
reconstruction imported from CT; FDG synthesis (1976) unlocking metabolic imaging;
NaI(Tl) and early BGO systems.}

\subsection{Key Milestones: Block Detectors, 3D Acquisition, Time-of-Flight}
\topics{the block detector (Casey \& Nutt 1986) and its cost significance; septa removal
and fully-3D acquisition; the first TOF wave (CsF/BaF\textsubscript{2}, 1980s) and why
it stalled; the LSO-enabled TOF revival (mid-2000s); whole-body FDG PET entering
routine oncology.}

\subsection{The Modern Era: Digital PET and Total-Body Systems}
\topics{SiPM-based ``digital'' PET (Philips Vereos, Siemens Vision, GE Discovery MI);
long axial field-of-view and total-body systems (uEXPLORER 2018, Biograph Quadra);
sensitivity gains and their price tags; current commercial landscape.}

\section{Physical Principles and Radiotracers}

\subsection{Positron Emission, Annihilation and Coincidence Detection}
\topics{$\beta^+$ decay and positron thermalisation; annihilation and back-to-back
511~keV photons; lines of response; coincidence timing windows; trues, randoms,
scattered and multiple coincidences.}

\subsection{Radionuclide Production: Cyclotrons and Generators}
\topics{cyclotron-produced isotopes (\textsuperscript{18}F, \textsuperscript{11}C,
\textsuperscript{15}O, \textsuperscript{13}N) and half-life logistics; generator systems
(\textsuperscript{68}Ga, \textsuperscript{82}Rb); distribution networks and cost;
implications for scanner siting.}

\subsection{Common Radiotracers and Their Applications}
\topics{\textsuperscript{18}F-FDG and glucose metabolism as the workhorse; PSMA ligands
in prostate cancer; amyloid and tau tracers in neurodegeneration; hypoxia and
proliferation tracers; theranostic pairs.}

\subsection{Physical Limits on Spatial Resolution: Positron Range and Acollinearity}
\topics{isotope-dependent positron range; $\sim$0.25\textdegree{} photon acollinearity
and its ring-diameter dependence; detector pitch contribution; the resolution budget
formula; why sub-mm whole-body PET is physics-limited.}

\subsection{Interaction of 511~keV Photons with Matter}
\topics{photoelectric absorption versus Compton scatter at 511~keV; linear attenuation
coefficients and photofraction by material; inter-crystal scatter; implications for
detector thickness and sensitivity.}

\section{Scintillators for PET}

\subsection{Scintillation Mechanisms and Key Performance Properties}
\topics{inorganic scintillation mechanism (band structure, activator dopants); light
yield, decay time, rise time; density and effective Z; emission spectrum matching to
photodetectors; energy resolution; afterglow; hygroscopicity and machinability;
figures of merit combining these.}

\subsection{Established Inorganic Crystals: NaI(Tl), BGO, LSO/LYSO and Alternatives}
\topics{property table across NaI(Tl), BGO, LSO:Ce, LYSO:Ce, GSO, BaF\textsubscript{2},
LaBr\textsubscript{3}; \textsuperscript{176}Lu intrinsic radioactivity; cost per unit
volume comparison; the BGO renaissance argument for non-TOF, high-sensitivity systems.}

\subsection{Emerging Materials: Transparent Ceramics, Nanocomposites and Plastics}
\topics{GAGG:Ce and garnet ceramics; transparent ceramic fabrication routes;
composites of inorganic nanoparticles in polymer (e.g.\ LaF\textsubscript{3}:Ce in
polystyrene) and their scatter and light-yield trade-offs; pure plastic scintillators
and the J-PET approach; manufacturability, moulding to arbitrary shapes, and cost.}

\subsection{Detector Segmentation: Pixellated Arrays versus Monolithic Slabs}
\topics{dicing, polishing and reflector assembly costs of pixellated arrays; dead space
and packing fraction; light sharing in monolithic slabs; flood-map calibration burden;
edge effects; intrinsic DOI information in monoliths; resolution below the SiPM pitch.}

\subsection{Depth-of-Interaction Encoding Approaches}
\topics{parallax error and radial elongation; phoswich and staggered-layer designs;
dual-ended readout; light-sharing-ratio methods in monoliths; cost and channel-count
implications of each approach.}

\section{Photodetectors and Readout Electronics}

\subsection{Photomultiplier Tubes and the Anger Logic Legacy}
\topics{PMT operation, gain and quantum efficiency; Anger centroid positioning; block
multiplexing ratios; why PMTs dominated for decades; magnetic field incompatibility.}

\subsection{Avalanche Photodiodes and Silicon Photomultipliers}
\topics{APDs as the transitional technology; SiPM/SPAD microcell operation; photon
detection efficiency, dark count rate, crosstalk, afterpulsing; temperature dependence
and bias control; analog versus digital SiPMs; MR compatibility; cost per channel
trends.}

\subsection{Readout ASICs and Channel Multiplexing}
\topics{TOFPET2/PETsys, Petiroc, STiC and similar ASICs; channel density and power;
resistive/capacitive multiplexing schemes; row--column and position-sensitive readout;
the channel count versus information trade-off.}

\subsection{Digitisation, Triggering and Coincidence Processing}
\topics{TDCs and time-over-threshold; FPGA-based coincidence engines; singles versus
coincidence data flow; data rates and online processing; open readout platforms used by
academic groups.}

\subsection{Timing Performance and Time-of-Flight Readout}
\topics{coincidence resolving time and its contributors (scintillator, SiPM,
electronics); the TOF SNR gain relation; current state of the art
($\sim$200--400~ps clinical); the 10~ps challenge; whether TOF is worth its cost in
cheap systems.}

\section{Scanner Architectures}

\subsection{The Conventional Cylindrical Clinical Scanner}
\topics{ring diameter and axial FOV conventions; crystal--block--module--ring
hierarchy; representative NEMA NU-2 figures of flagship systems; why the cylinder
became the default.}

\subsection{Preclinical and Small-Animal Systems}
\topics{NEMA NU-4 evaluation; sub-mm resolution requirements; small-bore geometry
effects (parallax, solid angle); representative commercial and academic systems;
preclinical as a proving ground for new detector technology.}

\subsection{Organ-Specific Systems: Brain, Breast and Prostate}
\topics{dedicated brain scanners (HRRT, NeuroEXPLORER, helmet/wearable concepts);
positron emission mammography; prostate probes; the proximity argument: solid angle and
resolution gains from conformal geometry; economic case for organ-specific systems.}

\subsection{Total-Body PET}
\topics{uEXPLORER and Biograph Quadra; order-of-magnitude sensitivity gains; dose
reduction, dynamic whole-body imaging and throughput economics; the inverse lesson for
low-cost design: sensitivity can be traded for crystal volume.}

\subsection{Unconventional and Experimental Geometries: Flat-Panel, Box and Enclosing Systems}
\topics{dual-panel coincidence benches as detector testbeds; box/enclosing geometries;
partial and sparse rings; plant PET, in-beam and proton-therapy range verification PET;
the monolithic annulus (AnnPET); reconstruction obstacles that non-cylindrical
geometries face with standard tools.}

\subsection{Case Studies of Research Scanner Development Programmes}
\topics{J-PET (plastic strips, axial arrangement, explicitly cost-driven); OpenPET and
open-source hardware efforts; MADPET and sparse-readout demonstrators; university-built
brain/breast prototypes; what these programmes reveal about the cost floor of
scanner construction and the path from bench to publishable system.}

\section{Image Reconstruction}

\subsection{Data Representations: Sinograms and List-Mode}
\topics{sinogram binning and michelograms; spatial quantisation error introduced by
binning; list-mode/event-by-event processing; compression versus information
preservation; where continuous LOR endpoint representations fit.}

\subsection{Analytical Reconstruction: Filtered Backprojection}
\topics{2D FBP and the central slice theorem; 3D reprojection; rebinning approaches
(SSRB, FORE); noise characteristics; why analytical methods persist (speed,
linearity, standardisation).}

\subsection{Iterative Reconstruction: MLEM, OSEM and Regularised Methods}
\topics{the Poisson likelihood model; MLEM derivation and convergence behaviour; OSEM
acceleration and subset artefacts; penalised/MAP methods (e.g.\ Q.Clear/BSREM);
stopping criteria and noise--resolution trade-offs.}

\subsection{System Matrix Modelling and Ray Tracing}
\topics{Siddon and Joseph projectors; on-the-fly computation versus precomputed/factorised
matrices; PSF/resolution modelling; GPU projector implementations; how the system
matrix encodes (and hard-codes) scanner geometry.}

\subsection{Data Corrections: Attenuation, Scatter, Randoms and Normalisation}
\topics{CT-based and MR-based attenuation correction; single scatter simulation;
delayed-window randoms estimation; component-based normalisation; dead time and decay
correction; motion correction in brief.}

\subsection{Time-of-Flight and Depth-of-Interaction in Reconstruction}
\topics{TOF kernels localising events along the LOR; improved conditioning and
convergence; incorporating DOI into the projector; joint benefit for unconventional
geometries.}

\subsection{Open-Source Reconstruction Packages}
\topics{STIR, CASToR, PyTomography, parallelproj, NiftyPET; scanner description
formats and their geometric assumptions; extensibility to novel geometries; the gap
between vendor and open-source toolchains.}

\subsection{Continuous-Coordinate and Geometry-Agnostic Reconstruction}
\topics{list-mode reconstruction with continuous detector coordinates; on-the-fly
Siddon on raw endpoint matrices; prior monolithic-annulus work that still binned to
sinograms; what has and has not been demonstrated; this as the primary methodological
gap the candidate's pilot work addresses.}

\section{Machine Learning in Modern PET}

\subsection{Neural Event Localisation in Monolithic Detectors}
\topics{CNN/MLP regression on SiPM light distributions; planar and DOI accuracy
reported in the literature; Compton-versus-photoelectric classification; timing
regression; uncertainty estimation; comparison with maximum-likelihood and
centre-of-gravity baselines.}

\subsection{Learned Reconstruction, Denoising and Image Enhancement}
\topics{direct DNN reconstruction (DeepPET); unrolled/model-based networks; U-Net and
diffusion denoisers; low-dose and fast-scan enhancement; hallucination risk, validation
standards and regulatory acceptance.}

\subsection{Monte Carlo Simulation as a Training Data Source}
\topics{Geant4/GATE and OpenGATE~10; optical photon transport fidelity; SiPM response
modelling; dataset generation at scale on HPC; simulation cost versus measured
calibration data.}

\subsection{The Simulation-to-Real Transfer Problem}
\topics{domain gap sources: surface finish, reflectors, SiPM non-uniformity, dark
counts, electronics noise; transfer learning and fine-tuning on measured floods;
calibration procedures; the sparse literature quantifying sim-to-real degradation ---
a candidate gap.}

\section{Contemporary Directions: How PET Is Being Improved}
\topics{synthesis section: reads across the pipeline sections above to identify where
the field is actively investing; each subsection cross-references the detailed
treatment earlier in the chapter and asks what the trend implies for cost.}

\subsection{Modern Approaches to Radiotracers}
\topics{theranostic pairs and the PSMA/DOTATATE boom; novel targets (FAPI, immuno-PET);
total-body PET enabling ultra-low-dose and dual-tracer protocols; decentralised/kit-based
tracer production and generator isotopes reducing dependence on cyclotron proximity.}

\subsection{Modern Approaches to Scintillators and Detectors}
\topics{the 10~ps timing challenge and fast light mechanisms (Cherenkov,
cross-luminescence, metascintillators and heterostructures); DOI-capable detector designs
going mainstream; monolithic detectors maturing from lab to product; cost-driven
material work: BGO revival for non-TOF systems, garnet ceramics, nanocomposites and
plastics.}

\subsection{Modern Approaches to Photodetectors and Readout}
\topics{SiPM cost and performance trajectory; digital SiPMs and 3D-integrated
photon-to-digital converters; high-density ASIC ecosystems; channel reduction via
multiplexing and on-detector intelligence; FPGA/firmware coincidence processing and
streaming readout.}

\subsection{Modern Approaches to Scanner Architecture}
\topics{long axial FOV / total-body as the headline trend and its cost countercurrent;
organ-specific and wearable systems; sparse, modular and upgradeable ring designs;
flat-panel and non-cylindrical concepts; PET/MR and multimodal integration.}

\subsection{Modern Approaches to Reconstruction and Data Processing}
\topics{TOF- and DOI-aware iterative reconstruction as standard; regularised methods in
clinical products; deep-learning reconstruction, denoising and low-dose imaging;
data-driven motion correction; geometry-agnostic and list-mode/continuous-coordinate
frameworks as the research frontier.}

\subsection{Implications for Low-Cost PET}
\topics{which of the above trends raise the performance ceiling versus lower the cost
floor; the recurring pattern of software substituting for hardware; where a
PhD-scale project can realistically contribute.}

\section{Clinical Relevance and Cost Drivers}

\subsection{Clinical Applications: Oncology, Neurology and Cardiology}
\topics{FDG staging, restaging and therapy response (Deauville criteria); dementia and
epilepsy workups; myocardial viability and perfusion; inflammation/infection imaging;
which applications tolerate lower performance and could be served by cheaper systems.}

\subsection{Sensitivity, Dose and Throughput Trade-offs}
\topics{NECR and count-rate behaviour; injected dose versus scan time versus image
quality; patient throughput as the economic driver of scan cost; where extra
sensitivity is best spent (dose, speed, or cheaper hardware).}

\subsection{Scanner Cost Structure: Crystals, Photodetectors and Assembly}
\topics{published bill-of-materials estimates; scintillator volume share of system
cost; SiPM and ASIC channel costs; assembly labour and calibration; gantry, cooling and
installation; service contracts and lifetime cost.}

\subsection{Accessibility and the Case for Low-Cost Architectures}
\topics{scanner density by region and the global access gap; PET in low- and
middle-income settings; mobile and modular concepts; what performance floor a
clinically useful cheap scanner must clear; how the candidate directions from
Chapter~\ref{ch:introduction} respond to this need.}

\section{Summary and Research Gap}
\topics{synthesis of the cost levers across materials, detectors, electronics, geometry
and algorithms; the specific under-explored combination this project targets;
restatement of candidate directions with a literature-informed assessment; criteria and
plan for converging on the PhD direction.}
